\section{Cuarta parte: hacia la asistencia de IA para matemáticas de nivel de investigación}

\subsection{Los retos particulares de la formalización de nivel de investigación}

La formalización de matemáticas de nivel de investigación (como el proyecto PFR de Terence Tao) implica:
\begin{enumerate}
    \item Crear un \textbf{plano} (Blueprint): descomponer la demostración en un grafo de dependencias
    \item Escribir una gran cantidad de \textbf{definiciones nuevas} y \textbf{lemas auxiliares}
    \item Ensamblar el teorema final y su demostración
\end{enumerate}

\begin{warningbox}{Limitaciones de las pruebas de referencia actuales}
Las pruebas de referencia más usadas (como mini-F2F) se basan en problemas de competencias matemáticas, que se caracterizan por ser autocontenidos y depender de resultados estándar. En cambio, los teoremas de los proyectos de investigación reales dependen de \textbf{definiciones y lemas nuevos propios del proyecto}; esta brecha hace que los modelos con buen desempeño en problemas de competencia puedan fallar gravemente en proyectos reales.
\end{warningbox}

\subsection{La prueba de referencia MiniCTX}

\begin{importantbox}{MiniCTX: una prueba de referencia de demostración de teoremas dependiente del contexto}
\begin{itemize}
    \item Extrae teoremas de prueba de proyectos reales de Lean (como Mathlib y PFR)
    \item Selecciona los teoremas por una fecha de corte, y puede actualizarse continuamente para mantenerse por delante de la fecha de corte del entrenamiento de los modelos
    \item Distingue entre dependencias dentro del mismo archivo y dependencias entre archivos
    \item Los experimentos muestran que: en problemas de competencia, el contexto casi no afecta el desempeño; en proyectos reales, \textbf{la falta de contexto provoca una caída abrupta del desempeño}
\end{itemize}
\end{importantbox}

\subsection{Herramienta práctica: LLMLean}

LLMLean es un complemento de Lean que puede conectarse a modelos locales (mediante Ollama o LM Studio) o a modelos de API (OpenAI, Claude, Gemini, etc.) y ofrece directamente en VS Code sugerencias de demostración verificadas por Lean.
